\scanpage{42}
% archive sheet number: 41
En particulier, si $G_0,G'_0$ satisfont aux conditions de la proposition, alors $\operatorname{Hom}(G_0,G'_0)\ne 0$, i.e. $G_0$ et $G'_0$ sont isogènes, ssi $\pi=\pi'$. \deleted{\unsure{Soit}} S'il en est ainsi, alors \deleted{\illegible{2}} on a
\[
\left\{\begin{aligned}
&\text{si }\nu\unsure{\leq|<}\nu'\quad\deleted{G_0=G'_0},\quad \operatorname{Hom}_A(\unsure{G_0|\Omega_0},G'_0)\xrightarrow{\sim}\operatorname{Hom}(\Omega',\Omega)\\
&\text{si }\nu>\nu'\qquad\operatorname{Hom}_A(G_0,G'_0)\simeq\pi\operatorname{Hom}(\Omega',\Omega)
\end{aligned}\right.
\]

\uline{Re NB} Revenons au cas général où on ne suppose pas néc. $\dim G_0=1$, \uuline{ni} $\varpi$ hauteur $=h$. Classifions à isogénie près \deleted{\unsure{les}} ces $G_0$, donc les $M_\nu$ $(0\leq\nu\leq f-1)$ deviennent des vectoriels \deleted{de \unsure{dim.}} \inserted{sur le} corps \deleted{\unsure{$A$}}$K$, isomorphes par les $\varphi_\nu$ $(0\leq\nu\leq f-2)$. Alors \deleted{\unsure{il est clair encore que}}
\[
M\simeq\deleted{\unsure{M}}\Omega\otimes_{\deleted{\unsure{\mathbb Z_p}}\mathbb Q_p}W\simeq\Omega\otimes_{\unsure{K_p|K}}\bigl(\deleted{\unsure{k}}K\deleted{\unsure{\otimes}}\otimes_{\mathbb Q_p}\underset{\inserted{\text{corps des fractions de } W}}{K_0}\bigr).
\]

et $F$ est déterminé par la connaissance de $\varphi_{f-1}$, qui est un endom de $M$, soit $\varpi$, de sorte qu'on a
\[
F_M^{\unsure{f|s}}=\varpi\otimes_{\mathbb Q_p}\mathrm{id}_{\deleted{\unsure{k}}K_0}=\varpi\otimes_K(K\otimes_{\mathbb Q_p}K_0)
\]
La condition \textquotedblleft d'effectivité\textquotedblright\ en termes de $\varpi$ est
\[
\left\{\begin{aligned}
&\varpi\ \text{laisse invariant un réseau}\\
&p^f\varpi^{-1}\ \text{laisse invariant un réseau}
\end{aligned}\right.
\]
